\documentclass[11pt]{article}
\usepackage[T1]{fontenc}
\usepackage[utf8]{inputenc}
\usepackage{lmodern,microtype,amsmath,amssymb,amsthm,geometry,booktabs}
\usepackage[colorlinks=true,linkcolor=blue,urlcolor=blue,citecolor=blue]{hyperref}
\geometry{margin=1in}
\setlength{\emergencystretch}{2em}
\newtheorem{theorem}{Theorem}
\newtheorem{lemma}[theorem]{Lemma}
\newtheorem{corollary}[theorem]{Corollary}
\newcommand{\R}{\mathbb R}
\newcommand{\E}{\mathbb E}
\newcommand{\norm}[1]{\left\lVert#1\right\rVert}
\title{Polynomial-depth Bayesian learning with general Lipschitz activations\\under balanced neuron-averaging stability}
\author{OpenAI Codex \and Kaarel H\"anni \and Dmitry Vaintrob}
\date{23 October 2026}
\hypersetup{pdftitle={Polynomial-depth Bayesian learning with general Lipschitz activations under balanced neuron-averaging stability},pdfauthor={OpenAI Codex; Kaarel Hanni; Dmitry Vaintrob}}
\begin{document}
\maketitle
\begin{abstract}
We give one polynomial sample-complexity argument for tanh, ReLU, exact GELU,
SiLU, sigmoid, softplus, and every fixed globally Lipschitz activation. No
zero-at-zero assumption is needed. The Gaussian prior is restricted to networks
whose clipped predictions survive averaging neurons within sufficiently few
balanced populations, with population count calibrated by readout energy.
A target implemented by a depth-$L$, width-$C$ network with weights bounded by
$B$ is learned with sample complexity polynomial jointly in $d,C,L,B$, the
activation's Lipschitz and offset bounds, and the statistical parameters.
One sample event works at every learner width $N\ge C$ and every prior scale
at least as strong as mean field, including all $\gamma\ge1$. The teacher
need not itself satisfy the restriction: a rescaled, cloned implementation
has a neighborhood of sufficient restricted prior mass. Its hidden states
may grow exponentially with depth; the resulting small parameter radius
costs only its logarithm. This is a theorem for a structurally restricted
prior, not for the unrestricted Gaussian learner or for internal-noise
stability alone. A separate appendix removes the zero-at-zero assumption
from the fixed-depth learning theorem for the original globally
margin-conditioned prior. The proofs are conventional and are not Lean formalized.
\end{abstract}

\section{The result and its setting}
The question is whether a single argument can give polynomial dependence on
depth for Bayesian networks with both bounded and unbounded activations.
We use a restriction expressing redundancy among hidden neurons. The same
balanced-population restriction was previously studied for tanh networks;
we show that its proof applies much more broadly. In particular, neither
positive homogeneity nor $\phi(0)=0$ is needed. The comparator here is an
actual bounded-weight neural teacher, so no Boolean-circuit compiler or
extra constant wire is hidden in the width threshold.

Fix $d,C,H\ge1$, $B,\lambda\ge1$, and $b\ge0$. Network depth is $L=H+1$.
Fix any activation $\phi:\R\to\R$ such that
\begin{equation}\label{eq:activation}
 |\phi(u)-\phi(v)|\le\lambda|u-v|,\qquad |\phi(0)|\le b.
\end{equation}
Inputs are $x\in\mathcal X=\{-1,1\}^d$, with $\bar x=(x,1)$.
A width-$N$ network has $H$ hidden layers and a linear readout:
\[
 h^1=\phi(W^1\bar x),\qquad h^s=\phi(W^sh^{s-1})\quad(2\le s\le H),
 \qquad z_w=v^\top h^H.
\]
Activation is coordinatewise; there are no separate internal biases.
The raw prior $\mu_{N,t}$ has independent centered Gaussian entries with
variances
\begin{equation}\label{eq:prior}
 t^2/d\text{ in }W^1,\qquad t^2/N\text{ in }W^s\ (s\ge2),
 \qquad t^2/N^2\text{ in }v,\qquad 0<t\le1.
\end{equation}
The scale may be prescribed as a function $t=t(d,L,N)$ of the
architecture. In particular, $t(N)=N^{(1-\gamma)/2}$ gives the exponent
$\gamma$ convention. The all-scale statements below are simultaneous
over their displayed range of $t$, so they include such varying choices
as well as a fixed scale; no constancy across widths is assumed. Define
\[
 E_v=N\norm v_2^2,\qquad T(u)=\max\{-1/2,\min\{1/2,u\}\}.
\]

A balanced partition into $q$ populations has sizes in
$\{\lfloor N/q\rfloor,\lceil N/q\rceil\}$. Let $P_s$ average the
coordinates within the populations at hidden layer $s$. For a tuple of
such partitions, all with the same $q$, define
\begin{equation}\label{eq:merged}
 \bar W^1=P_1W^1,\qquad \bar W^s=P_sW^sP_{s-1},\qquad \bar v=P_Hv.
\end{equation}
Let $\mathcal A_N^\phi$ be the event that some $q$ and balanced partitions
satisfy
\begin{equation}\label{eq:support}
 1\le q\le N,\qquad q\le1+E_v(w),\qquad
 \max_{x\in\mathcal X}|T(z_w(x))-T(z_{\bar w}(x))|\le1/16.
\end{equation}
This event depends on the activation and network, but not on the target
or training labels. It is Borel: there are finitely many witnesses at
fixed $N$, and the output maps are continuous. The condition uses absolute
readout energy; it is not invariant under arbitrary rescaling of the readout.

The prior is $\mu_{N,t}$ conditioned on $\mathcal A_N^\phi$. We also allow
the optional further restriction $|z_w(x)|>1/2$ at every input. Training
conditions on $f(X_i)z_w(X_i)>1/2$ for every example, where $X_i$ are iid
from any distribution $D$ on the cube. Write $Q_{S,N,t}$ for the resulting
posterior, with either convention. Test loss is
\[
 R_D(Q)=\E_{w\sim Q}\Pr_{X\sim D}\{f(X)z_w(X)\le1/4\}.
\]
This is a full posterior-average margin loss, which bounds sign error.

\begin{theorem}[One polynomial for all activations and all stronger scales]
\label{thm:poly}
There is a universal constant $K_0$ with the following property. Suppose
$f:\mathcal X\to\{-1,1\}$ has a width-$C$, $H$-hidden-layer implementation
using $\phi$, all of whose individual weights are bounded by $B$, and whose
output satisfies $f(x)z_*(x)>1/2$ at every input. Put
\[
 S_0=2+d+C+H+B+\lambda+b+\log(4/\varepsilon).
\]
For every $0<\varepsilon,\delta<1$, the sufficient condition
\begin{equation}\label{eq:poly-budget}
 n\ge\frac{K_0}{\varepsilon}
       \left[S_0^{13}+\log(1/\delta)\right]
\end{equation}
gives, with sample probability at least $1-\delta$,
$R_D(Q_{S,N,t})\le\varepsilon$ simultaneously for every $N\ge C$, every
$0<t\le1$, and both choices of the optional global-margin restriction.
All prior and training normalizers are positive. The sample polynomial
does not depend on $N$, $t$, or $\gamma$, and its coefficient and degree
are uniform over activations satisfying \eqref{eq:activation}.
\end{theorem}

The sample event may depend on the fixed activation and target. The
statement does not select an activation using the same data. The teacher
uses the learner's activation; no arbitrary Boolean circuit is identified
with a same-width implementation. Under the earlier complexity convention
$B=C$, the displayed envelope is still jointly polynomial in width and
depth. Its degree is deliberately conservative, not claimed optimal.

\begin{corollary}[Posterior confidence]
Apply Theorem~\ref{thm:poly} with accuracy $\varepsilon/3$. Then the posterior
probability of $f(x)z_w(x)>1/4$ is strictly greater than $2/3$ outside a set
of test mass at most $\varepsilon$. This also implies the same guarantee
for the correct sign. The conclusion allows arbitrary posterior mixtures.
\end{corollary}
\begin{proof}
If the correct-margin probability is at most $2/3$, the posterior error at
that input is at least $1/3$. Integrate this inequality under $D$.
\end{proof}

\section{Explicit budgets}
The following formulas make the polynomial statement concrete. Define
\begin{align*}
 R_*&=2+b+\lambda B(d+C+1),& M&=R_*^H,\\
 g_H&=2+\sqrt{2\log(2H)},&
 A_*&=1+\lambda(\sqrt2 CB+\sqrt C+g_H),\\
 V&=4\sqrt2 CB,&
 E&=\lambda(d+1+H\sqrt C)M A_*^{H-1},\\
 \rho&=\frac1{8[(V+1)E+M]},&
 B_0&=B^2[d(d+1)+(2H+30)C^2],\\
 q_0&=d+2+(H-1)C,&
 D_*&=B_0+q_0\log\frac{e\sqrt{2\pi}}{2\rho}+\log2.
\end{align*}
The potentially exponential quantities $M,E,1/\rho$ enter the prior
cost only through $\log(1/\rho)$. Set
\begin{align*}
 B_{r}&=\max\{8,8[D_*+\log(4/\varepsilon)]\},&
 Q&=\lceil1+B_{r}\rceil.
\end{align*}
\[
 J=8\left[D_*+\log(4/\varepsilon)+(H+1)\log Q
       +\tfrac{\log2}{2}\{d+2+(H-1)Q\}\right].
\]
\begin{align*}
 R&=1+\sqrt{QJ},&
 A_g&=2(1+b+\lambda)(d+1+Q)(R+1),\\
 \zeta&=[32(H+1)A_g^{H+1}]^{-1},&
 P_Q&=Q(d+2)+(H-1)Q^2.
\end{align*}

\begin{theorem}[Explicit sufficient sample size]\label{thm:explicit}
Every conclusion of Theorem~\ref{thm:poly} holds if
\begin{equation}\label{eq:explicit}
 n\ge\frac2\varepsilon
       \left[P_Q\log(3+4R/\zeta)+\log(1/\delta)\right].
\end{equation}
\end{theorem}

\section{A target neighborhood of polynomial logarithmic cost}
\begin{lemma}[Target mass]\label{lem:mass}
At every $N\ge C$ and $0<t\le1$, the raw Gaussian mass of networks in
$\mathcal A_N^\phi$ satisfying $f(x)z_w(x)>1$ everywhere is at least
$\exp(-ND_*/t^2)$.
\end{lemma}
\begin{proof}
Let $v$ be the teacher readout. Its norm is positive by the signed-margin
assumption. Replace it by
\[
 v'=4\max\{1,\norm v_2^{-1}\}v.
\]
Then $4\le\norm{v'}_2\le4B\sqrt C$, and the signed output is greater
than $2$. The multiplier itself can be large; the final coefficient norm,
which determines the Gaussian cost, is polynomially bounded.

Partition all $N$ neurons at each layer into $C$ balanced groups with
sizes $k_a$. Copy first-layer rows, put an internal teacher weight from
logical neuron $a$ to $b$ on each corresponding edge divided by $k_a$,
and divide each readout coefficient by its group's size. This preserves
every hidden value and the output exactly for every activation.
The cloned hidden state is constant within groups, and its coordinates
are bounded by $M$: the recurrence
$m_1\le b+\lambda B(d+1)$,
$m_s\le b+\lambda CBm_{s-1}$ gives $m_s\le R_*^s$.

For these cloned weights $w_*$,
\[
 \norm{W_*^s}_{\rm op}\le\sqrt2 CB\quad(s\ge2),\qquad
 8C\le N\norm{v_*}_2^2\le32C^2B^2,\qquad
 \sqrt N\norm{v_*}_2\le V.
\]
These follow from $C/2\le N/k_a\le2C$ and $k_b/k_a\le2$.
In independent Gaussian coordinates $\sqrt dW^1$, the coordinates of
each $\sqrt N W^s$ row along the $C$ normalized group indicators, and
$Nv$, constrain every coordinate within $\rho t$ of its center.
There are $Nq_0$ such coordinates, each with variance $t^2$, and the sum
of squared centers is at most $NB_0$.

Writing $P$ for the clone projection, an internal perturbed matrix is
$W^s=W_*^s+F^s+G^s$ with
\[
 F^s=F^sP,\quad \norm{F^s}_{\rm op}\le\rho t\sqrt C,
 \qquad G^sP=0.
\]
The free bulk $G^s=(t/\sqrt N)Z^s(I-P)$ is independent of the constrained
coordinates. The Gaussian operator-norm bound gives
$\Pr\{\norm{G^s}_{\rm op}>g_Ht\}\le1/(2H)$.
Thus all $H-1$ bulk bounds hold with conditional probability at least
$1/2$; for $H=1$ the assertion is immediate.

Put $e_s=\sup_x\norm{h_w^s(x)-h_*^s(x)}_2/\sqrt N$.
Lipschitz continuity, $G^sh_*^{s-1}=0$, and $t,\rho\le1$ give
\[
 e_1\le\lambda\rho(d+1),\qquad
 e_s\le A_*e_{s-1}+\lambda\rho\sqrt C M.
\]
Consequently $e_H\le\rho E$. The scaled readout error is at most $\rho$,
so the output error is at most
\[
 (V+1)e_H+\rho M\le1/8.
\]
The same estimate holds for the merged weights: the projections preserve
the constrained coordinate intervals, the centers, and kill the free bulk.
Both outputs therefore have signed value greater than $15/8$ everywhere.
Their clipped outputs agree exactly. Also
\[
 \sqrt{E_v(w)}\ge\sqrt{8C}-\rho\ge\sqrt{2C},
\]
since $\rho\le1/8$ and $C\ge1$. The clone partition thus witnesses
$C\le1+E_v(w)$ and \eqref{eq:support}. The optional global margin holds too.

For a scalar $Z\sim\mathcal N(0,t^2)$ and $\rho\le1$,
\[
 \Pr\{|Z-a|\le\rho t\}
 \ge\frac{2\rho}{e\sqrt{2\pi}}e^{-a^2/t^2}.
\]
Multiply these interval bounds and the independent bulk probability.
Using $t\le1$ and $N/t^2\ge1$ yields the claimed lower bound.
\end{proof}

\section{Gaussian control of every balanced partition}
For fixed partitions, the merged network is exactly a width-$q$ network
with the original activation and coefficients
\[
 U^1_{aj}=\frac1{k_{1,a}}\sum_{i\in a}W^1_{ij},\qquad
 U^s_{ab}=\frac1{k_{s,a}}\sum_{i\in a,j\in b}W^s_{ij},\qquad
 u_a=\sum_{j\in a}v_j.
\]
This identity holds for every activation, including nonzero offsets.
Its weighted energy is
\[
 \kappa_q=\frac dq\norm{U^1}_F^2+
             \sum_{s=2}^H\norm{U^s}_F^2+q\norm u_2^2.
\]
\begin{lemma}[Partition tail]\label{lem:partitions}
The raw probability that some balanced partitions with
$q\le\min\{N,Q\}$ have $\kappa_q>J$ is at most
\[
 Q\exp\left[HN\log Q+
 \tfrac N2\{d+2+(H-1)Q\}\log2-\frac{NJ}{8t^2}\right].
\]
\end{lemma}
\begin{proof}
For a fixed partition, compressed coordinates are independent centered
Gaussians of variances $t^2/(dk_{1,a})$,
$t^2k_{s-1,b}/(Nk_{s,a})$, and $t^2k_{H,a}/N^2$.
After the energy factors, each variance is at most $2t^2/N$ by balance.
There are $P_q=q(d+2)+(H-1)q^2$ coordinates, hence
\[
 \E e^{N\kappa_q/(8t^2)}\le2^{P_q/2}.
\]
There are at most $q^{HN}$ ordered partition tuples, and
$P_q/N\le d+2+(H-1)q$ for $q\le N$. Markov's inequality and the union
bound prove the assertion. The union is taken before choosing any
partition depending on the weights.
\end{proof}

\begin{lemma}[Uniform posterior truncation]\label{lem:truncation}
For every noiseless training sample, outside posterior mass at most
$\varepsilon/2$, a witnessing merged network has $q\le Q$ and
$\kappa_q\le J$, for every eligible width, scale and optional-margin choice.
\end{lemma}
\begin{proof}
The raw mass of each conditioning event is at least $e^{-ND_*/t^2}$ by
Lemma~\ref{lem:mass}. Since $NE_v/t^2$ is chi-square with $N$ degrees of
freedom, for $B_{r}\ge8$,
\[
 \mu_{N,t}\{E_v>B_{r}\}
 \le e^{-NB_{r}/(4t^2)}2^{N/2}
 \le e^{-NB_{r}/(8t^2)}.
\]
After normalization this is at most $\varepsilon/4$. Otherwise support
membership supplies $q\le1+B_{r}\le Q$.
Apply Lemma~\ref{lem:partitions}. Each positive exponent term, including
$\log Q$, is bounded by $N/t^2$ times its corresponding term in $J/8$.
The remaining raw probability is at most
$e^{-ND_*/t^2}(\varepsilon/4)^{N/t^2}$; after normalization it is at most
$\varepsilon/4$. Both estimates use only $N\ge1$ and $t\le1$.
\end{proof}

\section{Finite coding for activations with an offset}
A merged network with $q\le Q$ and energy at most $J$ has coefficients
bounded by $\sqrt{QJ}$. Pad to width $Q$, assigning zero to all edges
entering or leaving a dummy neuron and to its readout coefficient.
Dummy hidden outputs equal $\phi(0)$, but their outgoing weights are zero,
so padding preserves the output exactly.
Round coefficients on a fixed grid with mesh error at most $\zeta$,
preserving zero, and magnitude bounded by $R$. The class $\mathcal G$
of all grid networks has
\begin{equation}\label{eq:count}
 \log|\mathcal G|\le P_Q\log(3+4R/\zeta).
\end{equation}
Every network in this coding step still uses the one fixed activation.

The inequality $|\phi(u)|\le b+\lambda|u|$ gives hidden-coordinate
magnitudes at most $A_g^s$. If $e_s$ is the largest coordinate difference
between the padded and rounded networks, then
\[
 e_1\le\lambda(d+1)\zeta\le A_g\zeta,\qquad
 e_s\le\lambda QR e_{s-1}+\lambda Q\zeta A_g^{s-1}
      \le A_g e_{s-1}+\zeta A_g^s.
\]
Thus $e_s\le s\zeta A_g^s$, and the readout error is at most
$(H+1)\zeta A_g^{H+1}=1/32$. Merging and clipping give
\begin{equation}\label{eq:approx}
 |T(z_w(x))-T(g(x))|\le3/32
\end{equation}
for some $g\in\mathcal G$, uniformly over inputs. Only the logarithm of
the exponential precision enters \eqref{eq:count}.

\begin{proof}[Proof of Theorem~\ref{thm:explicit}]
For each grid predictor define
$F_g=\{x:f(x)T(g(x))\le11/32\}$.
If the original network fits the training margins and satisfies
\eqref{eq:approx}, then $f(X_i)T(g(X_i))\ge13/32$ for every training input,
so the sample avoids this event. Every test failure $f(x)z_w(x)\le1/4$
belongs to $F_g$. A fixed event of probability greater than
$\varepsilon/2$ is avoided by $n$ iid inputs with probability at most
$e^{-n\varepsilon/2}$. Union bounding over \eqref{eq:count} gives one
sample event of probability at least $1-\delta$ on which every retained
network has test loss at most $\varepsilon/2$.
The discarded posterior mass from Lemma~\ref{lem:truncation} adds at most
$\varepsilon/2$. The grid is independent of $N,t$ and the optional-margin
choice, proving simultaneity. Lemma~\ref{lem:mass} proves all normalizers
positive, even without a lower bound on sample size.
\end{proof}

\begin{proof}[Proof of Theorem~\ref{thm:poly}]
All constants in the following estimates are universal. For $S_0$ as in
the theorem, the target budgets satisfy
\[
 R_*\le2S_0^3,\quad A_*\le8S_0^3,\quad V\le6S_0^2,
 \quad \log(1/\rho)=O(H\log S_0),\quad D_*=O(S_0^5).
\]
The last bound follows from $B_0=O(S_0^5)$,
$q_0=O(S_0^2)$, and $H\log S_0=O(S_0^2)$.
Consequently
\[
 Q=O(S_0^5),\quad J=O(S_0^6),\quad R=O(S_0^6),\quad
 A_g=O(S_0^{12}),\quad P_Q=O(S_0^{11}).
\]
Finally $\log(3+4R/\zeta)=O(H\log S_0)=O(S_0^2)$, so the log-code
budget is $O(S_0^{13})$. Increase one universal constant to dominate the
explicit sample condition. No exponent depends on depth or activation.
\end{proof}

\section{Scope and interpretation}
The activation class includes all the examples below, with the indicated
conservative resource bounds.
\begin{center}
\begin{tabular}{lll}\toprule
Activation & Lipschitz bound $\lambda$ & Offset bound $b$\\\midrule
$\tanh(u)$, ReLU & $1$ & $0$\\
Exact GELU $u\Phi(u)$, SiLU $u/(1+e^{-u})$ & $2$ & $0$\\
Sigmoid $(1+e^{-u})^{-1}$ & $1$ & $1/2$\\
Softplus $\log(1+e^u)$ & $1$ & $\log2$\\\bottomrule
\end{tabular}
\end{center}
For a fixed activation its offset is a fixed constant. Removing
$\phi(0)=0$ therefore admits ordinary sigmoid and softplus without a
width penalty. Simply centering an activation is not an equality of the
original Bayesian learners: it creates hidden biases correlated with the
weights. Our argument analyzes the uncentered learner directly.

The prior restriction remains substantial. It requires balanced hidden
populations and prediction stability under their averaging, with the
count-energy relation \eqref{eq:support}. It is not implied here by
small internal-noise sensitivity, a bounded Jacobian, ordinary training,
or the mere existence of a small teacher. Its positive-mass teacher
construction establishes representability under the restriction, not
its prevalence among trained networks. The theorem does not retain an
additional arbitrary internal-noise restriction for general analog
teachers. Such an intersection needs its own target-mass argument.

This condition differs from the unsigned, arbitrary-population-size,
output-scale-invariant restriction in the ReLU/GELU polynomial-depth
paper. Balance bounds ratios of population sizes and lets us code the
original fixed activation directly; that is why no special perspective
identity for GELU, or homogeneity for ReLU, is needed. The sharper
quadratic Boolean-width results under those other conditions remain
separate. Nor does the present result settle unrestricted polynomial-depth
learning, efficient posterior inference, corrupted labels, or arbitrary
real-valued regression.

\paragraph{Relation to earlier work and verification.}
The balanced-population Gaussian tail and finite-code argument are inherited
from the collaborative tanh population paper, with contributions credited
there to ChatGPT-6 Astra and Claude. The new points here are the general
bounded-weight analog-teacher cylinder, its readout normalization, and
uniform coding for offset Lipschitz activations. These remove both the
digital-teacher restriction and the activation-specific parts of the proof.
The present extension has conventional proofs and AI-agent checks; the
earlier Lean certifications do not automatically certify it.

\appendix
\section{Nonzero activation offsets for the original prior at fixed depth}
\label{app:offset-original}

The restriction \eqref{eq:support} is unnecessary if depth is held fixed
and its polynomial exponent is allowed to depend on depth. This appendix
extends the unified fixed-depth theorem of~\cite{lipschitz}, whose stated
activation assumption was $\phi(0)=0$, to arbitrary bounded activation
offsets. It analyzes the original Gaussian learner directly. Centering
an activation and declaring the resulting network to have the same prior
would be incorrect, since it induces weight-dependent hidden biases.
The result below is a conventional proof extension; the earlier Lean
certification does not certify this new offset extension.

Keep the network and Gaussian law \eqref{eq:prior}, but impose
\emph{no internal-noise or neuron-averaging support restriction}. Instead,
let
\[
 U_{N,g}=\{w:g(x)z_w(x)>1/2\text{ for every }x\in\mathcal X\},
 \qquad U_N=\bigcup_g U_{N,g}.
\]
The prior is $\mu_{N,t}$ conditioned only on $U_N$. Given noiseless
examples of $f$, retain functions consistent with all training labels;
write the resulting function posterior as $\Pi_{S,N,t}$. Its risk is
\[
 R_D^{\rm sign}(\Pi)=\E_{g\sim\Pi}\Pr_{X\sim D}\{g(X)\ne f(X)\}.
\]
Unlike the main theorem's margin loss, this is sign error under a prior
already conditioned on the global absolute margin.

\begin{theorem}[Original-prior learning with bounded activation offsets]
\label{thm:offset-original}
For each fixed $L\ge2$, $\lambda\ge1$ and $b\ge0$, there is a polynomial
$\mathcal P_{L,\lambda,b}$ with the following property. Fix an activation
satisfying \eqref{eq:activation}, an input dimension $d\ge1$, and a
Boolean target admitting a depth-$L$, width-$C\ge1$ implementation with
individual weights bounded by $B\ge1$ and strict signed output margin
$1/2$ at every input. For every input law $D$ on the cube and every
$0<\varepsilon,\delta<1$,
\[
 n\ge\frac2\varepsilon\left[
 \mathcal P_{L,\lambda,b}(d,C,B,\log(1/\varepsilon))
 +\log(1/\delta)\right]
\]
implies, with training-sample probability at least $1-\delta$,
\[
 R_D^{\rm sign}(\Pi_{S,N,t})\le\varepsilon
 \qquad\text{simultaneously for every }N\ge C\text{ and }0<t\le1.
\]
All prior and posterior conditioning events have positive probability.
The polynomial is uniform over the activations satisfying the displayed
bounds. It includes every $\gamma\ge1$ and the strong-prior limit at each
fixed width, through $t=N^{(1-\gamma)/2}$ and $t\downarrow0$, respectively.
\end{theorem}

As before, applying the theorem at accuracy $\varepsilon/3$ gives
correct-label posterior probability strictly above $2/3$ outside test
mass at most $\varepsilon$. The teacher must use the learner's activation;
this does not assert a universal same-width circuit compiler for every
Lipschitz function. The polynomial degree in this appendix may grow with
$L$, in contrast with Theorem~\ref{thm:poly} under its additional support
restriction.

\subsection{A common capacity bound with offsets}
Put $H=L-1$ and retain the weighted energy
\[
 \kappa_N(w)=\frac dN\norm{W^1}_F^2+
 \sum_{s=2}^{H}\norm{W^s}_F^2+N\norm v_2^2.
\]
For counting only, allow each hidden neuron to use its own activation
with Lipschitz constant at most $\lambda$ and offset at most $b$.
Let $\mathcal F_b(K)$ contain all Boolean functions having a strict
margin-$1/2$ realization of energy at most $K$ in this auxiliary class,
at any finite equal hidden width. The actual learner still uses its
single fixed activation.

\begin{lemma}[Offset capacity]\label{lem:offset-capacity}
For $K\ge0$, define
\[
 \mathcal B_b(K)=
 \left\lceil\left[
 16H\lambda^H\sqrt{d+1}
 \left(\frac{(1+b^2)K+H}{L}\right)^{L/2}
 \right]^3\right\rceil\log(2^d+1).
\]
Then $\log|\mathcal F_b(K)|\le\mathcal B_b(K)$.
\end{lemma}
\begin{proof}
Write $a^s_j=\phi_{s,j}(0)$ and
$\psi_{s,j}=\phi_{s,j}-a^s_j$. For a width-$N$ realization put
$u^s=(h^s-a^s)/\sqrt N$. The state $(u^s,1)$ is computed by a
zero-at-zero network of the same depth, using upper-neuron activations
$q\mapsto\psi_{s,j}(\sqrt Nq)/\sqrt N$ and the identity activation for
its additional constant wire. Its first matrix consists of $W^1/\sqrt N$
and the row selecting the constant input. Its internal matrices and
readout are
\[
 \begin{pmatrix}W^s&W^sa^{s-1}/\sqrt N\\0&1\end{pmatrix},
 \qquad (\sqrt Nv^\top,\ v^\top a^H).
\]
These are exact identities, not approximations. Since
$\norm{a^s}_2\le b\sqrt N$, the sum of squared Frobenius norms of the
transformed layer matrices is at most $(1+b^2)\kappa_N+H$, using $d\ge1$.
Their product is therefore at most
$[((1+b^2)K+H)/L]^{L/2}$.

The zero-at-zero product-norm bound proved in~\cite{lipschitz} applies to
arbitrary neuronwise activations of Lipschitz constant $\lambda$:
on $q$ inputs of norm at most $R$, its empirical Rademacher complexity is
at most $8H\lambda^HPRq^{-1/3}$ when the product of layer norms is at
most $P$. Here $R=\sqrt{d+1}$. Shattering with strict margin $1/2$
forces that complexity to be at least $1/2$, giving the displayed VC bound.
The Sauer bound on the $2^d$-point cube proves the result.

The normalization by $\sqrt N$ is important: unnormalized centered
states would introduce columns $W^sa^{s-1}$ with an avoidable width
factor. The construction is used only to count functions; it is not an
identification of Gaussian priors.
\end{proof}

\subsection{Target probability and small-scale compression}
The target-cylinder argument of~\cite{lipschitz} requires only Lipschitz
bounds on differences once the teacher's states are bounded. Here the
teacher states satisfy
\[
 M_1\le b+\lambda B(d+1),\qquad
 M_s\le b+\lambda CBM_{s-1}.
\]
Thus a common bound $M$ is polynomial in $d,C,B$ at fixed $L,\lambda,b$.
Use balanced cloning into $C$ nonempty groups at every $N\ge C$;
this preserves the exact teacher function for any activation and uses
all neurons, so it does not assume zero padding is harmless.
Multiply the teacher readout by three.

For completeness, one may use
\[
 M=[1+b+\lambda B(d+C+1)]^H,\quad
 g_H=2+\sqrt{2\log(2H)},\quad
 F=1+\lambda(\sqrt2CB+\sqrt C+g_H),
\]
\[
 V=3\sqrt2CB,\qquad
 E=\lambda[(d+1)+H\sqrt C\,M]F^{H-1},\qquad
 \rho=[8((V+1)E+M)]^{-1}.
\]
Constrain first-layer, teacher-population, and readout coordinates
exactly as in the target-cylinder proof of~\cite{lipschitz}, to radius
$\rho t$ in coordinates with variance $t^2$. The free internal Gaussian
bulk annihilates teacher-group-constant states. With conditional
probability at least $1/2$, all its operator norms are at most $g_Ht$.
On that event the normalized state error obeys
\[
 e_1\le\lambda\rho(d+1),\qquad
 e_s\le Fe_{s-1}+\lambda\rho\sqrt C\,M,
 \qquad e_H\le\rho E.
\]
Hence output error is at most $(V+1)e_H+\rho M\le1/8$.
The draw has the teacher's signs with margin greater than $1/2$.

There are $ND_0$ constrained coordinates, where
$D_0=d+2+(H-1)C$, and their squared centers sum to at most $NB_0$, with
\[
 B_0=B^2[d(d+1)+(2H+36)C^2].
\]
The scalar interval estimate
\[
 \Pr\{|Z-a|\le\rho t\}\ge
 \frac{2\rho}{\sqrt{2\pi}}\exp(-a^2/t^2-\rho^2),
 \qquad Z\sim N(0,t^2),
\]
therefore gives
\begin{equation}\label{eq:offset-target}
 \mu_{N,t}(U_{N,f})\ge e^{-A_fN/(2t^2)},\qquad
 A_f=2\left[B_0+D_0\log\frac{e\sqrt{2\pi}}{2\rho}+\log2\right].
\end{equation}
The offset changes the state and radius budgets, but does not change the
Gaussian coordinates, independence, or cancellation of the free bulk.
All budgets have polynomial upper bounds at fixed depth.

The Gaussian spectral tail likewise makes no activation assumption.
On its usual good event the weighted first-layer and readout energies
are at most $K\ge1$, and each internal matrix is the sum of a retained
matrix of Frobenius norm at most $\sqrt K$ and a discarded matrix of
operator norm at most $a\le1$. Set
\[
 R=\lambda(\sqrt K+1)\ge2,\qquad U=2b+\lambda\sqrt{2K}.
\]
The original hidden RMS norms satisfy $M_s\le UR^{s-1}$ because
$M_1\le b+\lambda\sqrt{2K}$ and $M_s\le b+RM_{s-1}$.
Indeed, summing the additive offsets gives
$M_s\le R^{s-1}(b+\lambda\sqrt{2K})+b(R^{s-1}-1)/(R-1)
\le UR^{s-1}$.
For original versus truncated states,
\[
 e_1=0,\qquad
 e_s\le\lambda\sqrt K e_{s-1}+\lambda aM_{s-1}
 \le\lambda a(s-1)UR^{s-2}.
\]
Taking
\[
 a_K=\min\{1,[16\lambda(H+1)\sqrt K(1+U)R^H]^{-1}\}
\]
makes the readout discrepancy smaller than $1/8$. Doubling the readout
of the truncated network preserves signs and strict margin $1/2$, at
energy at most $(L+3)K$. The same Gaussian tails as in~\cite{lipschitz}
then imply
\begin{equation}\label{eq:offset-spectral}
 \mu_{N,t}\left(\bigcup_{g\notin\mathcal F_b((L+3)K)}U_{N,g}\right)
 \le Le^{-NK/(8t^2)}
\end{equation}
whenever $t\le a_K/8$ and $t^{-2}\ge8(d+1)$.
Here $a_K^{-1}$ is polynomially bounded at fixed depth. When $H=1$
there are no internal matrices to truncate and the discrepancy is zero.

\subsection{The mean-field endpoint with offsets}
This is the only part of the extension where an unmodified normalization
would fail. We give the replacement and then check every use of the
activation in the existing endpoint proof.

Let $Q$ be a positive semidefinite matrix on $m=2^d$ inputs, with diagonal
at most $M\ge1$, and write $\mathcal H_Q$ for its finite-dimensional
reproducing-kernel space. The Gaussian projection lemma of~\cite{lipschitz}
provides, for independent $Z_i\sim N(0,Q)$,
\[
 Z_i=a_i+u_i,\quad \norm{a_i}_\infty\le B_c:=\sqrt{2M\log(4m)},
 \quad \delta_i=\norm{u_i}_Q,
\]
with
\[
 \Pr\left\{N^{-1}\sum_i\delta_i^2>A\right\}\le e^{-Nq},
 \qquad A=4(q+\log2).
\]
For $Q=0$ take $a_i=u_i=0$. Keep its smooth cutoff $r_T$, which is one
on $\norm z_\infty\le T$ and zero outside $\norm z_\infty<\sqrt5T$.
Put $w_i=1-r_T(Z_i)^2$.

\begin{lemma}[Tail atoms with bounded offsets]\label{lem:offset-atoms}
Suppose $N^{-1}\sum_i\delta_i^2\le A$ and
$T\ge\max\{1,2B_c,b/\lambda\}$. The empirical activation covariance
differs in maximum norm from
\[
 \frac1N\sum_i r_T(Z_i)^2\phi(Z_i)\phi(Z_i)^\top
 +\frac1N\sum_i w_i\phi(u_i)\phi(u_i)^\top
\]
by at most $14\lambda^2B_cMA/T$. The second term has an exact representation
\[
 \sum_i c_i\psi_i(v_i)\psi_i(v_i)^\top,\qquad
 \sum_i c_i\le A+1,\quad \norm{v_i}_Q\le1,
 \quad \operatorname{Lip}(\psi_i)\le\lambda,\quad |\psi_i(0)|\le b.
\]
The next covariance diagonal is at most
$2(b+\lambda B_c)^2+2\lambda^2MA$.
\end{lemma}
\begin{proof}
When $w_i>0$, $\sqrt M\delta_i\ge T-B_c\ge T/2$. Consequently
\[
 N^{-1}\sum_iw_i\le4MA/T^2,\qquad
 N^{-1}\sum_iw_i\delta_i\le2\sqrt M A/T.
\]
Use $|\phi(Z_i)_x-\phi(u_i)_x|\le\lambda B_c$ and
$|\phi(u_i)_x|\le b+\lambda\sqrt M\delta_i$. The outer-product error is
bounded by
\[
 \frac{4\lambda^2B_cMA}{T}
 +\frac{8\lambda B_cbMA+4\lambda^2B_c^2MA}{T^2}
 \le\frac{14\lambda^2B_cMA}{T}.
\]
For the exact representation put
\[
 s_i=\sqrt{1+\delta_i^2},\qquad v_i=u_i/s_i,\qquad
 \psi_i(q)=\phi(s_iq)/s_i,\qquad c_i=w_is_i^2/N.
\]
Then $s_i\ge1$ proves all assertions, including
$\sum_i c_i\le1+A$. This normalization has no division by zero.
If $\delta_i=0$, then $w_i=0$; when $Q=0$, the whole tail is empty and
the bulk retains the actual constant feature $\phi(0)$.
Finally $|\phi(Z_i)_x|\le b+\lambda B_c+\lambda\sqrt M\delta_i$ gives
the diagonal bound.
\end{proof}

The stable RKHS certificate used in~\cite{lipschitz} is unchanged: for
$\norm v_Q\le1$ and $\eta>0$, it gives a vector $\beta$ such that
\[
 \norm{Q\beta-v}_\infty\le\eta,\qquad
 \norm\beta_1\le(4\eta)^{-1}.
\]
For a mixture of the atoms above, put $A'=A+1$ and
$R_A=b+\lambda\sqrt M$. Every atom has absolute values at most $R_A$.
For target error $0<\xi\le1$, choose
\[
 \eta=\frac{\xi}{64A'\lambda(R_A+\lambda+1)},\qquad
 n_A=\left\lceil128(A')^2R_A^4\xi^{-2}\log(4m^2)\right\rceil.
\]
Sample atoms with probabilities $c_i/A'$; fill the remaining probability
with vector zero \emph{and the identically zero activation}.
Hoeffding's inequality gives an empirical covariance approximation within
$\xi/4$. If $\norm{Q'-Q}_{\max}\le\eta^2$, the certificates satisfy
$\norm{Q'\beta-v}_\infty\le5\eta/4$. The activation error is at most
$5\lambda\eta/4$, and each outer-product error at most
$2R_A(5\lambda\eta/4)+(5\lambda\eta/4)^2$.
The choice of $\eta$ makes its total mixture contribution less than
$\xi/8$. Thus the stable finite-atom realization continues to work with
polynomial budgets. For a weighted feature realization
$Q'=\sum_\ell a_\ell h_\ell h_\ell^\top$, its incoming weights remain
\[
 W_{j\ell}=a_\ell\sum_y\beta_j(y)h_\ell(y),\qquad
 |W_{j\ell}|\le a_\ell\norm{h_\ell}_\infty/(4\eta).
\]
The dummy activation must be zero: using $\phi(0)$ for the dummy atom
would insert a spurious constant covariance.

\paragraph{The bounded bulk and its stability.}
Mollifying $\phi$ at scale $\tau\le1$ changes it by at most $\lambda\tau$,
retains first derivative bound $\lambda$, and gives second derivative
bound $O(\lambda/\tau)$. None of these facts uses $\phi(0)=0$.
For the cutoff features $b_x(z)=r_T(z)\phi_\tau(z_x)$ and the chosen
$T\ge b/\lambda$, the same entrywise derivative-sum bounds hold, with
larger universal constants:
\[
 B_0=O(\lambda T),\qquad B_1=O(\lambda),\qquad
 B_2=O\bigl(\lambda[\log(2m)/T+\tau^{-1}]\bigr).
\]
Indeed the additional offset terms in the first and second derivative
bounds are $O(b/T)$ and $O(b\log(2m)/T^2)$. The smoothing covariance
error is $O(\lambda^2T\tau)$.

The sparse exponential-tilt certificate and its Gaussian covariance
stability estimate in~\cite{lipschitz} apply to arbitrary bounded $C^2$
features. They do not assume a zero activation offset. For clarity,
with observables $F_{xy}=b_xb_y$, derivative budgets
$F_0=B_0^2$, $F_1=2B_0B_1$, $F_2=2B_1^2+2B_0B_2$, and tilt
$\norm\Gamma_{1,\rm ent}\le G$, their covariance map has modulus
\[
 C_G=\tfrac12\{F_2+2GF_1^2+2F_0(GF_2+G^2F_1^2)\}.
\]
It is polynomial in the indicated budgets, including for singular
covariances. The proof differentiates the normalized tilted expectation;
it does not pay an exponential reciprocal partition-function bound.

Given a desired one-step error $\alpha\le1$, take $q=4K$,
$A=4(q+\log2)$, and
\[
 T=\max\{1,2B_c,b/\lambda,256\lambda^2B_cMA/\alpha\}.
\]
Choose smoothing scale with reciprocal polynomial in $T,\alpha^{-1}$,
then $G=16q/\alpha$ and the sparse certificate size
$\max\{1,\lceil G^2F_0^2\rceil\}$. The original certificate proof has
exceptional mass at most $e^{-NK}$ above a polynomial width threshold.
Combine its error $\alpha/16$, smoothing error $\alpha/16$, and the tail
error from Lemma~\ref{lem:offset-atoms}; use the atom construction with
$\xi=\alpha/8$. Take the incoming covariance tolerance at most
$\min\{1,\eta^2,\alpha/[16(1+C_G)]\}$.

The bounded bulk is realized by finitely many tilted Gaussian rows,
evaluating the \emph{actual} activation $\phi$ and putting the cutoff
scalars into their covariance weights. Its feature bound is
$b+\lambda\sqrt5T=O(\lambda T)$. The Gaussian-coordinate intersection
argument of~\cite{lipschitz} bounds all sampled incoming coefficients
polynomially, since the square root of the tilt log-density bound enters
the coordinate cutoff. The expectation and sampling errors can each be
made at most $\alpha/16$ with polynomially many rows. All errors total
less than $\alpha$.

Thus the one-step construction has the same two exceptional events,
polynomial width and coefficient bounds, and a tolerance whose reciprocal
is polynomial in $K,M,\log(2m),\alpha^{-1}$ at fixed $\lambda,b$.
The physical hidden-value recursion becomes
$B_h^{\rm new}\le b+\lambda r B_W B_h$, where $r$ is the previous
finite feature count and $B_W$ the incoming weight bound. Its fixed-depth
iteration remains polynomial. Padding uses the identically zero
\emph{auxiliary} activation, which belongs to the bounded-offset class;
it does not introduce unused original-$\phi$ neurons with nonzero values.

\paragraph{Iteration and the final readout.}
Conditional on all previous layers, the next preactivation rows are
independent centered Gaussians with covariance $t^2Q_{s-1}$, where
\[
 Q_s(x,y)=N^{-1}\sum_i h_i^s(x)h_i^s(y),\qquad
 Q_0(x,y)=d^{-1}\bar x\cdot\bar y.
\]
These are raw Gram matrices, not centered hidden covariances; hidden
means therefore cause no change to this exact identity.
Iterate the diagonal bound in Lemma~\ref{lem:offset-atoms}, starting
from $\max Q_0(x,x)\le2$. Work backward from final covariance error
$1/(128K)$ using the reciprocal-polynomial one-step tolerances, then
reconstruct forward. Since $H$ is fixed and $\log m=d\log2$, all resulting
widths, coefficient bounds and physical feature values are polynomial in
$d,K$. The readout energy has law $t^2\chi_N^2/N$; except on mass
$e^{-NK}$ it is at most $8K$. The original output then has RKHS norm at
most $\sqrt{8K}$. Its unchanged certificate and the final covariance
approximation give output error at most $1/4$; doubling the reconstructed
readout recovers its signs with strict margin $1/2$.

All reconstructed activations have offset at most $b$ and Lipschitz
constant at most $\lambda$: bulk neurons use $\phi$, tail neurons use
the normalized $\psi_i$, and dummy/padding neurons use zero. Bounding
physical widths and weights bounds their weighted energy. Consequently
there are polynomials $P_{L,\lambda,b}$ and $J_{L,\lambda,b}$ such that
\begin{equation}\label{eq:offset-endpoint}
 N\ge P_{L,\lambda,b}(d,K)\quad\Longrightarrow\quad
 \mu_{N,t}\left(\bigcup_{g\notin\mathcal F_b(J_{L,\lambda,b}(d,K))}
 U_{N,g}\right)\le(2H+1)e^{-NK}.
\end{equation}
The conditional union bound accounts for two events per hidden layer
and one readout event. This holds uniformly for $0<t\le1$: covariance
multiplication by $t^2$ does not enlarge diagonal or approximation budgets,
bulk incoming weights acquire a factor $t$, atom incoming weights a factor
$t^2$, and readout tails only decrease. The offset itself is still bounded
by $b$.

\subsection{One finite class for every width and scale}
\begin{proof}[Proof of Theorem~\ref{thm:offset-original}]
Use $A_f$ from \eqref{eq:offset-target} and set
\begin{align*}
 K_s&=4A_f+8\log(2L/\varepsilon),&
 t_*&=\min\{a_{K_s}/8,[8(d+1)]^{-1/2}\},\\
 K_m&=2+A_f/t_*^2+\log(2(2H+1)/\varepsilon),&
 N_*&=\left\lceil\max\{1,P_{L,\lambda,b}(d,K_m)\}\right\rceil,\\
 D_*&=d+2+(L-2)N_*,&
 J_0&=2A_f+2D_*\log2+4\log(2/\varepsilon),\\
 J_*&=(L+3)K_s+J_{L,\lambda,b}(d,K_m)+J_0.
\end{align*}
Increase the polynomial envelopes to be nonnegative if needed. Every
quantity has a polynomial upper bound in $d,C,B,\log(1/\varepsilon)$
at fixed $L,\lambda,b$, including $t_*^{-1}$.

Let $\mathcal F_* =\mathcal F_b(J_*)$. For $t\le t_*$, divide the
spectral bound \eqref{eq:offset-spectral} by the target probability to get
\[
 \frac{\mu_{N,t}(\bigcup_{g\notin\mathcal F_*}U_{N,g})}
      {\mu_{N,t}(U_{N,f})}
 \le L\exp[-N(K_s/8-A_f/2)/t^2]\le\varepsilon/2.
\]
For $t\ge t_*$ and $N\ge N_*$, the endpoint bound
\eqref{eq:offset-endpoint} gives the same conclusion, since its ratio
is at most $(2H+1)\exp[-N(K_m-A_f/(2t^2))]$.
For $C\le N<N_*$, the parameter count is
$p=N(d+2)+(L-2)N^2\le ND_*$. Independently of the activation,
$\kappa_N=(t^2/N)\chi_p^2$ under the raw prior. Its exponential moment
at $N/(4t^2)$ is $2^{p/2}$, so
\[
 \frac{\mu_{N,t}\{\kappa_N>J_0\}}{\mu_{N,t}(U_{N,f})}
 \le \exp\left[-\frac N{t^2}(J_0/4-A_f/2)
                         +\frac p2\log2\right]\le\varepsilon/2.
\]
Every margin-separated draw of energy at most $J_0$ is itself in the
auxiliary bounded-offset class. Thus the same relative-mass bound holds
in every case.

Conditioning on the global margin cancels from probability ratios.
Noiseless training always leaves the target in the posterior denominator,
so posterior mass outside $\mathcal F_*$ is at most $\varepsilon/2$ on
\emph{every} training history. By Lemma~\ref{lem:offset-capacity},
$\log|\mathcal F_*|\le\mathcal B_b(J_*)$.
A finite-class union bound gives a sample event of probability at least
$1-\delta$ on which every consistent member has true error at most
$\varepsilon/2$, provided
\[
 n\ge\frac2\varepsilon[\mathcal B_b(J_*)+\log(1/\delta)].
\]
This event depends only on one finite function class and therefore works
for the entire family of widths and scales. Adding the discarded posterior
mass proves the risk guarantee. A polynomial upper envelope of
$\mathcal B_b(J_*)$ supplies $\mathcal P_{L,\lambda,b}$.
The target-cylinder lower bound proves all normalizers positive.
\end{proof}


\begin{thebibliography}{9}
\bibitem{populations} ChatGPT-6 Astra, Claude Fable 5 and Opus 5, K. H\"anni,
and D. Vaintrob. \emph{Quadratic-width learning bounds for Bayesian
networks with stable neuron populations}. 2026.
\href{https://github.com/kaarelh/AI-works/tree/main/neural-network-generalization/works-in-progress/improved-polynomials/neuron-averaging}{Paper and source}.
\bibitem{relugelu} ChatGPT-6 Astra, Claude Fable 5 and Opus 5, K. H\"anni,
and D. Vaintrob. \emph{Polynomial-depth learning for ReLU and GELU
NNBayes stable under neuron merging}. 2026.
\href{https://github.com/kaarelh/AI-works/tree/main/neural-network-generalization/works-in-progress/relu-gelu-polynomial-depth}{Paper and source}.
\bibitem{lipschitz} ChatGPT-6 Astra, Claude Fable 5 and Opus 5, K. H\"anni,
and D. Vaintrob. \emph{Width-independent learning bounds for ReLU, GELU,
and other Lipschitz activations}. 2026.
\href{https://github.com/kaarelh/AI-works/tree/main/neural-network-generalization/works-in-progress/unbounded-activations}{Paper and source}.
\end{thebibliography}
\end{document}
